\documentclass[10pt,a4paper]{amsart}
\usepackage[margin=19mm]{geometry}
\usepackage{amsmath,amssymb,mathtools}
\usepackage[scr=rsfs,cal=euler]{mathalfa}
\usepackage{xfrac}
\usepackage[unicode,hidelinks,bookmarks=false]{hyperref}
\newcommand{\ie}{i.e.\ }
\newcommand{\cf}{cf.\ }
\newcommand{\resp}{respectively\ }
\newcommand{\Subsection}{\subsection}
\providecommand{\nobreakdash}{\nobreak\hbox{-}}
\setlength{\emergencystretch}{2em}
\multlinegap0pt
\usepackage{amssymb}
\newtheorem{proposition}{Proposition}
\newcommand{\E}{\mathbb{E}}
\newcommand{\KL}{\mathrm{KL}}
\newcommand{\given}{\,\vert\,}
\title{Marginally Useful:\\ An Information-Gap Identity in Split Conformal Prediction}
\author{Peter Cotton}
\date{Source: arXiv:2608.07479v2 (CC BY 4.0)}
\begin{document}
\maketitle

\noindent\emph{Source and licence.} Marginally Useful:\\ An Information-Gap Identity in Split Conformal Prediction. Version arXiv:2608.07479v2 (CC BY 4.0), distributed by arXiv under the Creative Commons Attribution 4.0 International licence, http://creativecommons.org/licenses/by/4.0/, as recorded in the arXiv licence metadata for this exact version and shown on the abstract page https://arxiv.org/abs/2608.07479v2. Full licence notice: \texttt{licenses/arxiv-2608.07479-CC-BY-4.0.txt}.\par\smallskip

\noindent\textit{Editorial scope.} Counted unit: the article's proposition prop:gap (source0.tex lines 174--184). The article's complete original proof of it is delivered with it (source0.tex lines 186--202, one \texttt{proof} environment of 17 lines). No mathematics is written, repaired or re-derived: the statement and the proof are the article's own lines, in the article's own notation, and no retained line is substituted or re-flowed.  No further source-local context is needed: every cross-reference of every retained line resolves inside the two delivered spans.  Nothing was omitted from any retained span, so the package carries no omission record.  No retained line contains a citation, so the wrapper carries no bibliography and no bibliography entry.  No retained line contains a figure environment, an image-inclusion command or drawing code, so no image is delivered and the package declares no asset dependency.  Editorial formatting only, in addition: an AMS \texttt{amsart} preamble that restates the article's own declarations and macros used by the retained lines, namely the environments \texttt{proposition} on separate counters, together with the standard packages those lines need.  The retained environments print in this document as Proposition 1 in order of appearance, whereas the article prints the same items with the numbering of its own full document.  The complete native source of the article as distributed in the arXiv e-print for this exact version is bundled unchanged as \texttt{sources/207090/source0.tex}.
\begin{proposition}[The residual-information gap]\label{prop:gap}
With $I(R;X)=\E_X\,\KL(r(\cdot\given X)\,\|\,\bar r)\ge 0$ the mutual information between
the residual and the input, and assuming the relevant densities and information quantities
are finite,
\begin{equation}\label{eq:gap}
  \E[\log q^\star(Y\given X)]-\E[\log \pi_X(Y)] \;=\; I(R;X).
\end{equation}
Among all single-shape forecasters $y\mapsto h(y-\hat\mu(x))$, $\E[\log h(R)]$ is maximized
at $h=\bar r$, and the signed CPS attains this optimum. Hence no recalibration that ignores
$X$ can reduce $I(R;X)$. Reducing it requires conditioning on $X$.
\end{proposition}

\begin{proof}
Let $\mathcal{R}(h)$ be the expected log-score regret of the single-shape forecaster
$y\mapsto h(y-\hat\mu(x))$ relative to the oracle. The oracle scores $Y$ by
$\log q^\star(Y\given X)=\log r(R\given X)$ and this forecaster by $\log h(R)$. The joint
$P_{X,R}$ has density $p(x)\,r(\rho\given x)$ and the product reference $P_X\otimes H$ has
density $p(x)\,h(\rho)$, so
\begin{equation}\label{eq:master}
  \mathcal{R}(h)=\E\!\left[\log\frac{r(R\given X)}{h(R)}\right]
      =\KL\!\big(P_{X,R}\,\big\|\,P_X\otimes H\big)
      = I(R;X)+\KL(\bar r\,\|\,h),
\end{equation}
the last step the chain rule for relative entropy along a product reference, the $R$-marginal
of $P_{X,R}$ being $\bar r$. Taking $h=\bar r$ kills the second term and gives
\eqref{eq:gap}. That second term is $\ge 0$ and is minimized at $h=\bar r$ by Gibbs, so no
$X$-blind shape improves on $I(R;X)$, and $I(R;X)=\KL(P_{X,R}\,\|\,P_X P_R)$ is reduced only
by conditioning on $X$.
\end{proof}

\end{document}
